%% kapa: exam=1 round=21 year=2010 period=1 form=A subject=civil_law subject_ko=민법 seq=2 number=2
%% source: Q-Net 감정평가사 기출문제 (cst00309) artlSeq=5024181 file=extracted/1008/5024181/2010년 제21회 감정평가사 1차시험 기출문제/1교시.hwp
%% conversion: kordoc 4.24.51 -> tools/gen1.py (auto); review=proofread-sonnet; flags=none
\begin{question}{2}
하자있는 의사표시에 관한 설명으로 옳지 않은 것은? (다툼이 있으면 판례에 의함)
\begin{choices}
  \item 표의자가 사기ㆍ강박을 이유로 의사표시를 취소하지 않는 한 그 법률행위는 유효하다.
  \item 상대방이 없는 의사표시에 관하여 제3자가 사기나 강박을 한 때에는 특별한 사정이 없는 한 취소할 수 있다.
  \item 토지거래허가를 받지 않아 유동적 무효 상태인 매매계약이더라도, 사기ㆍ강박에 의한 의사표시임을 이유로 그 계약을 취소할 수 있다.
  \item 기망행위를 한 자가 상대방의 대리인인 경우, 상대방이 그 사실을 알았거나 알 수 있었을 경우에 한하여 그 의사표시를 취소할 수 있다.
  \item 중대한 고지의무 위반은 부작위에 의한 기망행위에 해당될 수 있기 때문에, 표의자는 의사표시를 취소하거나 기망행위로 인한 손해배상을 청구할 수 있다.
\end{choices}
\end{question}
\answer{4}
